\documentclass[12pt, a4paper]{ctexart}

%==================== 导言区：宏包加载 ====================
\usepackage{amsmath, amssymb, amsthm} % 数学公式与符号
\usepackage{geometry}                 % 页面设置
\usepackage{xcolor}                   % 颜色支持
\usepackage{fancyhdr}                 % 页眉页脚
\usepackage{enumitem}                 % 列表定制
\usepackage{tcolorbox}                %以此美化文本框
\usepackage{titlesec}                 % 标题格式调整

%==================== 页面布局设置 ====================
\geometry{left=2.5cm, right=2.5cm, top=2.5cm, bottom=2.5cm}

%==================== 自定义视觉样式 ====================

% 1. 定义分隔线样式
\newcommand{\TopSeparator}{
    \par\noindent\rule{\textwidth}{1.5pt}\vspace{0.5em} % 粗实线 ━━━━
}
\newcommand{\AnalysisSeparator}{
    \par\noindent\rule{\textwidth}{0.5pt}\vspace{0.2em}\hrule\vspace{0.5em} % 双线效果 ════
}

% 2. 定义红色解析环境
%在这个环境内的所有文字和公式都会自动变成红色
\newenvironment{RedAnalysis}{
    \par
    \color{red} % 设置后续字体为红色
}{
    \par
    \color{black} % 结束时恢复黑色
}

% 3. 标题格式微调
\titleformat{\section}{\large\bfseries\heiti}{}{0em}{}[\vspace{-0.5em}]
\titleformat{\subsection}{\bfseries\kaishu}{}{0em}{}

\newcommand{\choicesFour}[4]{
    \begin{enumerate}[label=\Alph*., itemsep=0pt, parsep=0pt, topsep=5pt, labelsep=0.5em]
        \begin{minipage}{0.22\linewidth} \item #1 \end{minipage}
        \begin{minipage}{0.22\linewidth} \item #2 \end{minipage}
        \begin{minipage}{0.22\linewidth} \item #3 \end{minipage}
        \begin{minipage}{0.22\linewidth} \item #4 \end{minipage}
    \end{enumerate}
}
\newcommand{\choicesTwo}[4]{
    \begin{enumerate}[label=\Alph*., itemsep=0pt, parsep=0pt, topsep=5pt, labelsep=0.5em]
        \begin{minipage}{0.45\linewidth} \item #1 \end{minipage}
        \begin{minipage}{0.45\linewidth} \item #2 \end{minipage}
        \begin{minipage}{0.45\linewidth} \item #3 \end{minipage}
        \begin{minipage}{0.45\linewidth} \item #4 \end{minipage}
    \end{enumerate}
}
\newcommand{\choices}[4]{
    \begin{enumerate}[label=\Alph*., itemsep=0pt, parsep=0pt, topsep=5pt, labelsep=0.5em]
        \item #1
        \item #2
        \item #3
        \item #4
    \end{enumerate}
}

%==================== 文档开始 ====================
\begin{document}

% 标题信息
\title{\textbf{第三章 连续（基础）解析讲义}}
\author{数学习题讲义}
\date{\today}
\maketitle

% --- 题目 1 ---
\TopSeparator
\section*{题目1}

\textbf{【题目重现】} \\
函数 $ f(x) $ 在 $\mathbb{R}$ 上连续， $x=2$ 是 $ f(x)=\frac{x^{2}-4}{x-2} $ 的 \underline{\hspace{1cm}}。
\choicesFour{可去间断点}{跳跃间断点}{无穷间断点}{连续点}

\AnalysisSeparator

\begin{RedAnalysis}

\subsection*{一、逐步解析}

\textbf{第一步：问题分析} \\
判断函数的间断点类型。函数在 $x=2$ 处未定义，但题目称函数在 $\mathbb{R}$ 上连续？
题目表述可能是：“若某函数在 $x \ne 2$ 时表达式为...”，或者单纯问这个表达式定义的函数在 $x=2$ 处的特性，并补充定义使其连续。
按通常理解，考查该表达式在 $x=2$ 处的极限是否存在。

\textbf{第二步：思路构建} \\
计算 $\lim_{x\to 2} f(x)$。
- 若极限存在，为可去间断点。
- 若极限不存在且左右极限不相等，为跳跃。
- 若极限无穷大，为无穷间断点。

\textbf{第三步：详细步骤} \\
1. 计算极限：
   \[ \lim_{x\to 2} \frac{x^2-4}{x-2} = \lim_{x\to 2} \frac{(x-2)(x+2)}{x-2} = \lim_{x\to 2} (x+2) = 4 \]
2. 极限存在。
3. 结论：$x=2$ 是可去间断点（可以通过补充定义 $f(2)=4$ 使其连续）。

\textbf{第四步：最终答案} \\
A

\subsection*{二、核心公式}
1. \textbf{可去间断点}：$\lim_{x\to x_0} f(x) = A$ (存在)，但 $f(x_0) \ne A$ 或未定义。

\subsection*{三、考点分析}
\begin{itemize}
    \item \textbf{知识点归类}：间断点分类。
\end{itemize}

\end{RedAnalysis}

\vspace{2em}

% --- 题目 2 ---
\TopSeparator
\section*{题目2}

\textbf{【题目重现】} \\
$ x \rightarrow 0^{+} $ 时，下列函数 \underline{\hspace{1cm}} 是无穷小量。
\choicesFour{$ e^{\frac{1}{x}} $}{$ x\sin\frac{1}{x} $}{$ \ln x $}{$ \frac{1}{x}\sin x $}

\AnalysisSeparator

\begin{RedAnalysis}

\subsection*{一、逐步解析}

\textbf{第一步：问题分析} \\
判断当 $x \to 0^+$ 时，哪个函数的极限为 0。

\textbf{第二步：思路构建} \\
逐一计算极限。

\textbf{第三步：详细步骤} \\
1. A选项：$\lim_{x\to 0^+} e^{1/x} = e^{+\infty} = +\infty$（无穷大）。
2. B选项：$x$ 是无穷小，$\sin(1/x)$ 有界。无穷小 $\times$ 有界 = 无穷小。$\to 0$。正确。
3. C选项：$\lim_{x\to 0^+} \ln x = -\infty$（无穷大）。
4. D选项：$\lim_{x\to 0^+} \frac{\sin x}{x} = 1$（非零常数）。

\textbf{第四步：最终答案} \\
B

\subsection*{二、核心公式}
1. \textbf{无穷小性质}：无穷小 $\times$ 有界 = 无穷小。

\subsection*{三、考点分析}
\begin{itemize}
    \item \textbf{知识点归类}：无穷小量的判定。
\end{itemize}

\end{RedAnalysis}

\vspace{2em}

% --- 题目 3 ---
\TopSeparator
\section*{题目3}

\textbf{【题目重现】} \\
当 $ x \rightarrow 0 $ 时，与 $ x + x^{3} $ 等价的无穷小量为 \underline{\hspace{1cm}}。
\choicesFour{$ \frac{1}{x} $}{1}{$ x $}{$ x^{2} $}

\AnalysisSeparator

\begin{RedAnalysis}

\subsection*{一、逐步解析}

\textbf{第一步：问题分析} \\
寻找主部。

\textbf{第二步：思路构建} \\
当 $x \to 0$ 时，低阶无穷小起决定作用。

\textbf{第三步：详细步骤} \\
$x + x^3 = x(1 + x^2)$。
因为 $1+x^2 \to 1$，所以极限比值为 1。
$x+x^3 \sim x$。

\textbf{第四步：最终答案} \\
C

\subsection*{二、核心公式}
略

\subsection*{三、考点分析}
\begin{itemize}
    \item \textbf{知识点归类}：无穷小比阶。
\end{itemize}

\end{RedAnalysis}

\vspace{2em}

% --- 题目 4 ---
\TopSeparator
\section*{题目4}

\textbf{【题目重现】} \\
函数 $ f(x)=\lim_{n\to\infty}\frac{1+x}{1+x^{2n}} $ 的间断点情况是 \underline{\hspace{1cm}}。
\choicesTwo
    {不存在间断点}
    {存在间断点 $ x = 1 $}
    {存在间断点 $ x = 0 $}
    {存在间断点 $ x = -1 $}

\AnalysisSeparator

\begin{RedAnalysis}

\subsection*{一、逐步解析}

\textbf{第一步：问题分析} \\
需要先求出 $f(x)$ 的解析式，讨论 $x$ 在不同范围内 $x^{2n}$ 的极限。

\textbf{第二步：思路构建} \\
分 $|x| < 1$, $|x| > 1$, $|x| = 1$ 讨论。

\textbf{第三步：详细步骤} \\
1. 求解析式：
   - 当 $|x| < 1$ 时，$x^{2n} \to 0$， $f(x) = \frac{1+x}{1} = 1+x$。
   - 当 $|x| > 1$ 时，$x^{2n} \to +\infty$， $f(x) = 0$。
   - 当 $x = 1$ 时，$1^{2n}=1$， $f(1) = \frac{2}{2} = 1$。
   - 当 $x = -1$ 时，$(-1)^{2n}=1$， $f(-1) = \frac{0}{2} = 0$。
2. 分析间断点：
   - 在 $x=1$ 处：
     左极限 $\lim_{x\to 1^-} (1+x) = 2$。
     右极限 $\lim_{x\to 1^+} 0 = 0$。
     左右不相等，为跳跃间断点。
   - 在 $x=-1$ 处：
     左极限 $\lim_{x\to -1^-} 0 = 0$。
     右极限 $\lim_{x\to -1^+} (1+x) = 0$。
     $f(-1)=0$。
     三者相等，连续。
3. 结论：只是在 $x=1$ 处间断。

\textbf{第四步：最终答案} \\
B

\subsection*{二、核心公式}
1. \textbf{幂函数极限}：$\lim_{n\to\infty} x^n$ 依赖于 $|x|$ 与 1 的关系。

\subsection*{三、考点分析}
\begin{itemize}
    \item \textbf{知识点归类}：极限函数与连续性。
    \item \textbf{易错点提示}：容易想当然认为 $x=-1$ 也是间断点，必须经过计算验证。
\end{itemize}

\end{RedAnalysis}

\vspace{2em}

% --- 题目 5 ---
\TopSeparator
\section*{题目5}

\textbf{【题目重现】} \\
设 $ x_{n}=\frac{n}{2}[1+(-1)^{n}] $ ，则 \underline{\hspace{1cm}}。
\choicesTwo
    {$ \{x_{n}\} $ 有界}
    {$ \{x_{n}\} $ 无界}
    {$ \{x_{n}\} $ 单调增加}
    {$ n\to\infty $ 时， $ \{x_{n}\} $ 为无穷大}

\AnalysisSeparator

\begin{RedAnalysis}

\subsection*{一、逐步解析}

\textbf{第一步：问题分析} \\
观察数列通项。

\textbf{第二步：思路构建} \\
列举几项或分奇偶性讨论。

\textbf{第三步：详细步骤} \\
- 当 $n$ 为偶数时，$(-1)^n = 1$，$x_n = \frac{n}{2} \cdot 2 = n \to \infty$。
- 当 $n$ 为奇数时，$(-1)^n = -1$，$x_n = 0$。
数列包含趋于无穷的子列，故无界。
但因有为0的子列，整体并不趋于无穷大（无穷大要求所有项绝对值都无限增大）。
显然不单调。

\textbf{第四步：最终答案} \\
B

\subsection*{二、核心公式}
略

\subsection*{三、考点分析}
\begin{itemize}
    \item \textbf{知识点归类}：数列性质判别。
\end{itemize}

\end{RedAnalysis}

\vspace{2em}

% --- 题目 6 ---
\TopSeparator
\section*{题目6}

\textbf{【题目重现】} \\
$ x = 0 $ 是函数 $ f(x) = \frac{\tan x}{x} $ 的第 \underline{\hspace{1cm}} 类间断点。

\AnalysisSeparator

\begin{RedAnalysis}

\subsection*{一、逐步解析}

\textbf{第一步：问题分析} \\
计算该点的极限。

\textbf{第二步：思路构建} \\
$\lim_{x\to 0} \frac{\tan x}{x} = 1$。
极限存在，但函数无定义。

\textbf{第三步：详细步骤} \\
极限存在，故为第一类间断点中的可去间断点。
属于第一类。

\textbf{第四步：最终答案} \\
一（或可去）

\subsection*{二、核心公式}
略

\subsection*{三、考点分析}
略

\end{RedAnalysis}

\vspace{2em}

% --- 题目 7 ---
\TopSeparator
\section*{题目7}

\textbf{【题目重现】} \\
$ f(x)=\begin{cases}x\sin\frac{1}{x}, & x\neq0,\\ k, & x=0\end{cases} $ 在 $x=0$ 处连续,则常数 $k=$ \underline{\hspace{2cm}}。

\AnalysisSeparator

\begin{RedAnalysis}

\subsection*{一、逐步解析}

\textbf{第一步：问题分析} \\
连续的充要条件是 $\lim_{x\to x_0} f(x) = f(x_0)$。

\textbf{第二步：思路构建} \\
计算 $\lim_{x\to 0} x\sin\frac{1}{x}$。

\textbf{第三步：详细步骤} \\
$\lim_{x\to 0} x\sin\frac{1}{x} = 0$（无穷小 $\times$ 有界）。
故须 $k=0$。

\textbf{第四步：最终答案} \\
0

\subsection*{二、核心公式}
略

\subsection*{三、考点分析}
略

\end{RedAnalysis}

\vspace{2em}

% --- 题目 8 ---
\TopSeparator
\section*{题目8}

\textbf{【题目重现】} \\
已知当 $ x \to 0 $ 时， $ (1 + ax^{2})^{3} - 1 $ 与 $ 1 - \cos x $ 是等价无穷小，则常数 $ a = $ \underline{\hspace{2cm}}。

\AnalysisSeparator

\begin{RedAnalysis}

\subsection*{一、逐步解析}

\textbf{第一步：问题分析} \\
利用常见等价无穷小建立等式。

\textbf{第二步：思路构建} \\
1. $(1+u)^3 - 1 \sim 3u$。
2. $1 - \cos x \sim \frac{1}{2}x^2$。

\textbf{第三步：详细步骤} \\
$(1 + ax^{2})^{3} - 1 \sim 3(ax^2) = 3ax^2$。
由题意：$3ax^2 = \frac{1}{2}x^2$。
$3a = \frac{1}{2} \Rightarrow a = \frac{1}{6}$。

\textbf{第四步：最终答案} \\
$1/6$

\subsection*{二、核心公式}
1. $(1+x)^\alpha - 1 \sim \alpha x$。
2. $1-\cos x \sim \frac{1}{2}x^2$。

\subsection*{三、考点分析}
略

\end{RedAnalysis}

\vspace{2em}

% --- 题目 9 ---
\TopSeparator
\section*{题目9}

\textbf{【题目重现】} \\
$ f(x)=\begin{cases}(1-x)^{\frac{1}{x}}, & x<0,\\ 2^{x}+k, & x\geq0\end{cases} $ 在 $ x=0 $ 处连续,则常数 $ k= $ \underline{\hspace{2cm}}。

\AnalysisSeparator

\begin{RedAnalysis}

\subsection*{一、逐步解析}

\textbf{第一步：问题分析} \\
连续 $\Rightarrow$ 左极限 = 右极限 = 函数值。

\textbf{第二步：思路构建} \\
分别求左极限和右极限。

\textbf{第三步：详细步骤} \\
1. 左极限：$\lim_{x\to 0^-} (1-x)^{\frac{1}{x}}$。
   令 $t = -x$，则 $\lim_{t\to 0^+} (1+t)^{-1/t} = [(1+t)^{1/t}]^{-1} = e^{-1} = \frac{1}{e}$。
2. 右极限：$\lim_{x\to 0^+} (2^x + k) = 1 + k$。
3. 令 $1+k = \frac{1}{e} \Rightarrow k = \frac{1}{e} - 1$。

\textbf{第四步：最终答案} \\
$\frac{1}{e}-1$

\subsection*{二、核心公式}
1. $\lim (1+x)^{1/x} = e$。

\subsection*{三、考点分析}
略

\end{RedAnalysis}

\vspace{2em}

% --- 题目 10 ---
\TopSeparator
\section*{题目10}

\textbf{【题目重现】} \\
函数 $ y=\frac{x}{\tan x} $ 的间断点为 \underline{\hspace{2cm}}。

\AnalysisSeparator

\begin{RedAnalysis}

\subsection*{一、逐步解析}

\textbf{第一步：问题分析} \\
寻找函数无定义的点。
1. 分母 $\tan x = 0$。
2. $\tan x$ 自身无定义的点。

\textbf{第二步：思路构建} \\
1. $\tan x = 0 \Rightarrow x = k\pi$。
2. $\tan x$ 无定义 $\Rightarrow \cos x = 0 \Rightarrow x = k\pi + \frac{\pi}{2}$。

\textbf{第三步：详细步骤} \\
在这些点处函数均无定义，均为间断点。
题目若填空，涵盖所有点集：$x = \frac{k\pi}{2}$ ($k \in \mathbb{Z}$)。

\textbf{第四步：最终答案} \\
$x = \frac{k\pi}{2}, k\in\mathbb{Z}$

\subsection*{二、核心公式}
略

\subsection*{三、考点分析}
略

\end{RedAnalysis}

\vspace{2em}

% --- 题目 11 ---
\TopSeparator
\section*{题目11}

\textbf{【题目重现】} \\
已知函数 $ f(x)=\begin{cases}(1+ax)^{\frac{1}{x}}, & x>0\\ 2b-e, & x=0\\ b+\ln(1+x^{2}), & x<0\end{cases} $ 在 $x=0$ 处连续，求实数 $a$ 和 $b$ 的值。

\AnalysisSeparator

\begin{RedAnalysis}

\subsection*{一、逐步解析}

\textbf{第一步：问题分析} \\
连续性条件的综合应用。

\textbf{第二步：思路构建} \\
$\lim_{x\to 0^+} f(x) = \lim_{x\to 0^-} f(x) = f(0)$。

\textbf{第三步：详细步骤} \\
1. 右极限：$\lim_{x\to 0^+} (1+ax)^{1/x} = e^a$。
2. 左极限：$\lim_{x\to 0^-} (b + \ln(1+x^2)) = b + 0 = b$。
3. 函数值：$f(0) = 2b - e$。
4. 联立方程：
   $e^a = b$
   $b = 2b - e \Rightarrow b = e$
5. 代入得 $e^a = e \Rightarrow a = 1$。

\textbf{第四步：最终答案} \\
$a=1, b=e$

\subsection*{二、核心公式}
略

\subsection*{三、考点分析}
略

\end{RedAnalysis}

\vspace{2em}

% --- 题目 12 ---
\TopSeparator
\section*{题目12}

\textbf{【题目重现】} \\
当 $ a, b $ 为何值时， $ \lim_{x \to 0} \frac{\sin x}{a - e^{x}} (b - \cos x) = 2 $。

\AnalysisSeparator

\begin{RedAnalysis}

\subsection*{一、逐步解析}

\textbf{第一步：问题分析} \\
根据极限值反推参数。

\textbf{第二步：思路构建} \\
分析分母极限为 0 的情况（若分母不为0则极限为0不等于2，除非分子有异）。

\textbf{第三步：详细步骤} \\
1. 分母 $a - e^x \to a-1$。若 $a \ne 1$，分母非零，分子 $\to 0 \cdot (b-1) = 0$。极限为0。与题意矛盾。
   故必有 $a - 1 = 0 \Rightarrow a = 1$。
2. 当 $a=1$ 时，分母 $1-e^x \sim -x$。
   原式 $= \lim \frac{\sin x}{-x}(b-\cos x) = \lim (-1) \cdot (b-\cos 0) = -(b-1) = 1-b$。
3. 由题意 $1-b = 2 \Rightarrow b = -1$。

\textbf{第四步：最终答案} \\
$a=1, b=-1$

\subsection*{二、核心公式}
略

\subsection*{三、考点分析}
略

\end{RedAnalysis}

\vspace{2em}

% --- 题目 13 ---
\TopSeparator
\section*{题目13}

\textbf{【题目重现】} \\
$ f(x) = \begin{cases} \frac{\cos x}{x+2}, & x \geq 0 \\ \frac{\sqrt{a} - \sqrt{a - x}}{x}, & x < 0 \quad (a > 0) \end{cases} $
\begin{enumerate}
    \item 当 $a$ 为何值时， $x=0$ 是 $ f(x) $ 的连续点？
    \item 当 $a$ 为何值时， $x=0$ 是 $ f(x) $ 的间断点？是什么类型的间断点？
\end{enumerate}

\AnalysisSeparator

\begin{RedAnalysis}

\subsection*{一、逐步解析}

\textbf{第一步：问题分析} \\
比较左右极限和函数值。

\textbf{第二步：思路构建} \\
计算右极限（即函数值）和左极限。

\textbf{第三步：详细步骤} \\
1. 右极限与函数值：
   $f(0) = \frac{\cos 0}{0+2} = \frac{1}{2}$。
   $\lim_{x\to 0^+} f(x) = \frac{1}{2}$。
2. 左极限：
   $\lim_{x\to 0^-} \frac{\sqrt{a}-\sqrt{a-x}}{x} = \lim \frac{a-(a-x)}{x(\sqrt{a}+\sqrt{a-x})} = \lim \frac{1}{\sqrt{a}+\sqrt{a-x}} = \frac{1}{2\sqrt{a}}$。
3. 讨论：
   - (1) 连续 $\iff$ 左=右：$\frac{1}{2\sqrt{a}} = \frac{1}{2} \Rightarrow \sqrt{a}=1 \Rightarrow a=1$。
   - (2) 间断 $\iff a \ne 1$。此时左右极限均存在但不相等，为跳跃间断点。

\textbf{第四步：最终答案} \\
(1) $a=1$; (2) $a \ne 1$ 时为跳跃间断点。

\subsection*{二、核心公式}
略

\subsection*{三、考点分析}
略

\end{RedAnalysis}

\vspace{2em}

% --- 题目 14 ---
\TopSeparator
\section*{题目14}

\textbf{【题目重现】} \\
讨论极限 $ \lim_{x \to 0} \frac{1 - 2^{\frac{1}{x}}}{1 + 2^{\frac{1}{x}}} $ 是否存在？

\AnalysisSeparator

\begin{RedAnalysis}

\subsection*{一、逐步解析}

\textbf{第一步：问题分析} \\
讨论左右极限。

\textbf{第二步：思路构建} \\
$x$ 趋于 $0^+$ 和 $0^-$ 时 $1/x$ 趋向不同的无穷大。

\textbf{第三步：详细步骤} \\
1. $x \to 0^+$ 时，$\frac{1}{x} \to +\infty$，$2^{1/x} \to +\infty$。
   极限 $= \lim \frac{2^{-1/x} - 1}{2^{-1/x} + 1} = -1$。
2. $x \to 0^-$ 时，$\frac{1}{x} \to -\infty$，$2^{1/x} \to 0$。
   极限 $= \frac{1-0}{1+0} = 1$。
3. 左右不相等，不存在。

\textbf{第四步：最终答案} \\
不存在

\subsection*{二、核心公式}
略

\subsection*{三、考点分析}
略

\end{RedAnalysis}

\vspace{2em}

% --- 题目 15 ---
\TopSeparator
\section*{题目15}

\textbf{【题目重现】} \\
研究函数 $ f(x)=\lim_{n\to\infty}\frac{1-x^{2n}}{1+x^{2n}}x $ 的连续性，并画出图形。

\AnalysisSeparator

\begin{RedAnalysis}

\subsection*{一、逐步解析}

\textbf{第一步：问题分析} \\
先求极限函数 $f(x)$。

\textbf{第二步：思路构建} \\
分段讨论 $|x|<1$, $|x|>1$, $|x|=1$。

\textbf{第三步：详细步骤} \\
1. 表达式：
   - $|x| < 1$: $x^{2n} \to 0 \Rightarrow f(x) = x$。
   - $|x| > 1$: $x^{2n} \to \infty \Rightarrow f(x) = -x$（分子分母同除以 $x^{2n}$）。
   - $|x| = 1$: $x^{2n} = 1 \Rightarrow f(x) = 0$。
2. 连续性：
   在 $(-\infty, -1), (-1, 1), (1, \infty)$ 上显然连续。
   在 $x=1$：左极限 $1$, 右极限 $-1$, $f(1)=0$。不连续。
   在 $x=-1$：左极限 $-(-1)=1$, 右极限 $-1$, $f(-1)=0$。不连续。
   即在 $x=\pm 1$ 处间断。
3. 图形：
   中间一段 $y=x$，两头 $y=-x$，连接点处取 0。

\textbf{第四步：最终答案} \\
在 $x=\pm 1$ 处不连续，其余处连续。

\subsection*{二、核心公式}
略

\subsection*{三、考点分析}
\begin{itemize}
    \item \textbf{知识点归类}：函数列极限与分段函数连续性。
\end{itemize}

\end{RedAnalysis}

\end{document}
